\chapter{\emph{k-hop reachability} en detalle}
\label{cap:khop-detalle}

\begin{intuicion}
k-hop reachability es el índice más complejo del motor: bitmap
de tipos alcanzables, layout dense+sparse para ahorrar memoria,
con extensión temporal (\texttt{min\_first}, \texttt{max\_first})
para podar también por ventana. Capítulo~\ref{cap:prune-khop} lo
introdujo como prune; este capítulo entra a la implementación
concreta: cómo se computa el bitmap con BFS reverse multi-source,
por qué el rank array es \texttt{u32} con sentinel
\texttt{NO\_RANK}, y un ejemplo numérico de la poda temporal en
acción.
\end{intuicion}

\section{La estructura}

\begin{lstlisting}[language=Rust]
const MAX_DEPTH: usize = 4;
const NO_RANK: u32 = u32::MAX;
const NLF_TYPE_COUNT: usize = 9;  // = |L_V|

pub struct KHopReach {
    pub bits: Vec<u16>,                      // [|V|], dense
    rank: Vec<u32>,                          // [|V|], dense; NO_RANK si outdeg=0
    min_first: RwLock<IndexBacking<[i64; NLF_TYPE_COUNT]>>,  // sparse
    max_first: RwLock<IndexBacking<[i64; NLF_TYPE_COUNT]>>,  // sparse
}
\end{lstlisting}

\noindent
Cuatro arrays:

\begin{itemize}
    \item \texttt{bits}: 2 bytes por vértice. \(\bits[v]\) tiene
    bit \(L\) prendido sii existe walk de \(v\) a algún vértice
    de tipo \(L\) en \(\leq \texttt{MAX\_DEPTH}\) saltos.
    \item \texttt{rank}: 4 bytes por vértice.
    \texttt{rank[v] = NO\_RANK} si \(v\) no tiene aristas salientes
    (no aporta info temporal); si no, indexa una fila en las
    matrices sparse de \texttt{min/max\_first}.
    \item \texttt{min\_first}: matriz sparse de \texttt{[i64; 9]};
    fila \(r\) es válida sólo si algún vértice \(v\) tiene
    \texttt{rank[v] = r}. Bytes 72 por fila.
    \item \texttt{max\_first}: igual.
\end{itemize}

\subsection{Por qué dense+sparse}

El layout naive ---dense para los cuatro arrays--- pesaría
\(|\V| \cdot (2 + 4 + 72 + 72)\) = \(|\V| \cdot 150\) B. Para
\(|\V| = 10^8\): 15 GB. Inviable.

La observación clave: \texttt{min\_first} y \texttt{max\_first}
sólo aportan info para vértices con \emph{aristas salientes}.
Vértices ``hojas'' (\(\outdeg = 0\)) tienen
\(\mathrm{min/max\_first}\) sin sentido. En grafos típicos el
\(\sim 70\%\) de los vértices son hojas en algún sentido.

El layout sparse: \texttt{rank} es dense, pero las matrices grandes
sólo guardan filas para los vértices ``ranked''. Memoria:
\(|\V| \cdot 6 + |\V_{\text{ranked}}| \cdot 144\). Para \(|\V| =
10^8\) y \(|\V_{\text{ranked}}| = 3 \cdot 10^7\): \(\sim 4.9\) GB.
\(3\times\) menos.

\section{El builder: BFS reverse multi-source}

\begin{algorithm}[H]
\caption{\texttt{KHopReach::build} (versión simplificada)}
\KwIn{grafo \(G = (\V, \E)\)}
\KwOut{\texttt{KHopReach}}
\(\bits[v] \leftarrow 0\) para todo \(v\)\;
\(\mathrm{min\_first}[v][L] \leftarrow +\infty\), \(\mathrm{max\_first}[v][L] \leftarrow -\infty\) para todo \((v, L)\)\;
construir \texttt{reverse adjacency} \(N^{-}\)\;
\BlankLine
\For{cada tipo \(L \in \Lc_V\)}{
    \(\mathrm{frontier}_0 \leftarrow \{v : \mathrm{type}(v) = L\}\)\;
    \For{\(d = 0, 1, \dots, \texttt{MAX\_DEPTH}\)}{
        \For{cada \(w \in \mathrm{frontier}_d\)}{
            \(\bits[w] \mathrel{|}\!= (1 \ll L)\)\;
        }
        \(\mathrm{frontier}_{d+1} \leftarrow \emptyset\)\;
        \For{cada \(w \in \mathrm{frontier}_d\)}{
            \For{cada \(u \in N^{-}(w)\)}{
                \(\tau \leftarrow \mathrm{timestamp\,de\,arista}\;u \to w\)\;
                \(\mathrm{min\_first}[u][L] \leftarrow \min(\mathrm{min\_first}[u][L], \tau)\)\;
                \(\mathrm{max\_first}[u][L] \leftarrow \max(\mathrm{max\_first}[u][L], \tau)\)\;
                agregar \(u\) a \(\mathrm{frontier}_{d+1}\) si no estaba\;
            }
        }
    }
}
asignar ranks: \(r \leftarrow 0\); para cada \(v\) con \(\outdeg(v) > 0\), \(\mathrm{rank}[v] \leftarrow r\), incrementar \(r\)\;
compactar matrices sparse a \(r\) filas\;
\Return\;
\end{algorithm}

\noindent
Costo: \(O(|\Lc_V| \cdot \texttt{MAX\_DEPTH} \cdot |\E|)\). Para
\(|\Lc_V| = 9\), \(\texttt{MAX\_DEPTH} = 4\), \(|\E| = 10^8\):
\(\sim 3.6 \cdot 10^{10}\) operaciones; $\sim$1.5 minutos en
hardware típico.

\subsection{Por qué \emph{reverse} BFS}

La pregunta natural: ``existe walk de \(v\) a tipo \(L\)'' es
direccional desde \(v\). Una BFS \emph{forward} desde \(v\)
respondería ``cuáles tipos alcanza \(v\)''. Pero queremos
responder para \emph{todos} los \(v\) sin pagar \(|\V|\) BFS
forward.

La trampa: corremos BFS \emph{reverse} desde cada conjunto
\(\{v : \mathrm{type}(v) = L\}\). Eso visita ``cuáles vértices
pueden alcanzar tipo \(L\)'' en una sola pasada, marcando el bit
\(L\) en cada uno. Para \(|\Lc_V|\) tipos, \(|\Lc_V|\) BFS
reverse total ---fixed cost.

\section{Lookup y prune}

\begin{lstlisting}[language=Rust]
impl KHopReach {
    pub fn can_reach(&self, v: StrId, target_type: u8) -> bool {
        self.bits[v.index()] & (1 << target_type) != 0
    }

    pub fn within_window(
        &self,
        v: StrId,
        target_type: u8,
        t_lo: i64,
        t_hi: i64,
    ) -> bool {
        if !self.can_reach(v, target_type) {
            return false;
        }
        let rank = self.rank[v.index()];
        if rank == NO_RANK {
            return false;  // outdeg = 0
        }
        let mn_lock = self.min_first.read();
        let mx_lock = self.max_first.read();
        let mn = mn_lock.get(rank as usize).map(|r| r[target_type as usize]);
        let mx = mx_lock.get(rank as usize).map(|r| r[target_type as usize]);
        // Si no hay info temporal, no se puede descartar: aceptar.
        let (mn, mx) = match (mn, mx) {
            (Some(m), Some(M)) => (m, M),
            _ => return true,
        };
        // Sound: descartar si min_first >= t_hi o max_first < t_lo
        !(mn >= t_hi || mx < t_lo)
    }
}
\end{lstlisting}

\noindent
\texttt{can\_reach}: bitwise \texttt{and}, una sola instrucción.
\texttt{within\_window}: una vez chequeado el bit, dos lookups
sparse + dos comparaciones \texttt{i64}. \(O(1)\) total.

\section{Ejemplo numérico}

\begin{ejemplo}[Prune temporal en acción]
Grafo de 5 vértices. Aristas (con timestamps):
\begin{itemize}
    \item \(u \to h_1\) en \(\tau = 100\)
    \item \(u \to h_2\) en \(\tau = 500\)
    \item \(h_1 \to ip_1\) en \(\tau = 200\)
    \item \(h_2 \to ip_2\) en \(\tau = 600\)
\end{itemize}

\noindent
Tipos: \(\mathrm{type}(u) = \texttt{User}\) (tag 0),
\(\mathrm{type}(h_i) = \texttt{Host}\) (tag 1),
\(\mathrm{type}(ip_i) = \texttt{IP}\) (tag 3).

Build de KHopReach (con \texttt{MAX\_DEPTH} = 2):

\begin{itemize}
    \item BFS reverse desde \(\{ip_1, ip_2\}\) marcando tag IP:
    \begin{itemize}
        \item Distancia 0: \(\{ip_1, ip_2\}\). \(\bits[ip_1]\) y
        \(\bits[ip_2]\) prenden bit 3.
        \item Distancia 1: predecesores son \(\{h_1, h_2\}\).
        \(\bits[h_i] \mathrel{|}\!= 8\) (bit 3 prendido).
        \(\mathrm{min\_first}[h_1][\texttt{IP}] = 200\),
        \(\mathrm{min\_first}[h_2][\texttt{IP}] = 600\). \(\mathrm{max\_first}\)
        igual (sólo una arista cada uno).
        \item Distancia 2: predecesores de \(h_1\) y \(h_2\) =
        \(\{u\}\). \(\bits[u] \mathrel{|}\!= 8\).
        \(\mathrm{min\_first}[u][\texttt{IP}] = \min(100, 500) = 100\).
        \(\mathrm{max\_first}[u][\texttt{IP}] = \max(100, 500) = 500\).
    \end{itemize}
\end{itemize}

\noindent
Hipótesis: \texttt{User -[Auth]-> Host -[Connect]-> IP WITHIN [t\_lo=300, t\_hi=400)}.

Prune temporal sobre \(u\) como start:
\begin{itemize}
    \item \texttt{can\_reach(u, IP)}: bit 3 prendido $\Rightarrow$
    true.
    \item \(\mathrm{min\_first}[u][\texttt{IP}] = 100 \geq 400\)?
    No (100 < 400).
    \item \(\mathrm{max\_first}[u][\texttt{IP}] = 500 < 300\)? No
    (500 \(\geq\) 300).
    \item Prune \emph{no descarta}: \(u\) podría tener un walk en
    la ventana.
\end{itemize}

\noindent
El DFS verifica concretamente: ¿hay arista \(u \to ?\) con
\(\tau \in [300, 400)\)? Las dos aristas salientes son \(\tau = 100\)
y \(\tau = 500\), ninguna en la ventana. Falso positivo del
prune; el DFS lo descarta tras verificar.

\textbf{El prune es conservador}: no falsea negativos, pero acepta
algunos falsos positivos. El DFS los limpia.

Si la ventana fuera \([t_{\mathrm{lo}}=600, t_{\mathrm{hi}}=700)\):
\begin{itemize}
    \item \(\mathrm{max\_first}[u][\texttt{IP}] = 500 < 600\) $\Rightarrow$
    prune \emph{descarta} \(u\) sin verificar. Ahorro real.
\end{itemize}
\end{ejemplo}

\section{El \texttt{RwLock} en min/max\_first}

¿Por qué min/max\_first están envueltos en \texttt{RwLock}? Porque
el plan futuro (capítulo~\ref{cap:delta-wcoj}) introduce delta
update incremental sobre k-hop. Las mutaciones son raras pero
existen; el \texttt{RwLock} permite lecturas concurrentes (matchers)
con escrituras exclusivas (delta update).

Hoy, sin delta update, los \texttt{RwLock} están en modo
read-only puro: el cost del lock acquire es $\sim$5 ns por
lookup, despreciable.

\section{Memoria efectiva}

Para grafo de 28 GB EVTX (\(|\V| = 10^8\), 30\%
\emph{outgoing-having} \(\Rightarrow\) \(|\V_{\text{ranked}}| = 3 \cdot
10^7\)):

\begin{center}
\renewcommand{\arraystretch}{1.2}
\begin{tabular}{lr}
\toprule
\textbf{Componente} & \textbf{Memoria} \\
\midrule
\texttt{bits: Vec<u16>} & 200 MB \\
\texttt{rank: Vec<u32>} & 400 MB \\
\texttt{min\_first: IndexBacking<[i64;9]>} (sparse) & 2.16 GB \\
\texttt{max\_first: IndexBacking<[i64;9]>} (sparse) & 2.16 GB \\
\midrule
\textbf{Total} & \(\sim 4.9\) GB \\
\bottomrule
\end{tabular}
\end{center}

\noindent
Las dos sparse matrices spillean a mmap (1 GB budget cada una).
Los dos arrays dense (bits, rank, suma 600 MB) viven en heap.

\section{Resumen}

\begin{itemize}
    \item k-hop reachability: \texttt{bits[v]} (u16, dense) +
    \texttt{rank[v]} (u32, dense) + \texttt{min/max\_first} (i64
    arrays, sparse).
    \item Layout dense+sparse: \(\sim 3\times\) menos memoria que
    dense puro.
    \item Build: BFS reverse multi-source, una pasada por tipo.
    \(O(|\Lc_V| \cdot \texttt{MAX\_DEPTH} \cdot |\E|)\).
    \item Lookup: bitwise para can\_reach, lookup sparse + dos
    comparaciones para temporal.
    \item \texttt{RwLock} en min/max\_first: hoy read-only, futuro
    para delta update (capítulo~\ref{cap:delta-wcoj}).
\end{itemize}

\begin{ejercicios}
\begin{enumerate}
    \item \starred{} Para un grafo \(|\V| = 5 \cdot 10^7\) con
    50\% \emph{outgoing-having}, calculá la memoria total de
    \texttt{KHopReach}.

    \item ¿Por qué BFS \emph{reverse} y no \emph{forward}? Pensá:
    para responder ``puede \(v\) alcanzar tipo \(L\)?'' para todos
    los \(v\) y todos los \(L\), ¿cuál es más eficiente?

    \item \starred{} Diseñá un caso adversarial donde el prune
    temporal NO descarta nada aunque debería. ¿Es esto un bug o un
    trade-off de la sobre-aproximación?

    \item Si \texttt{MAX\_DEPTH = 6} en lugar de 4, ¿cuánto crece
    el tiempo de build? ¿Y la memoria?

    \item Imaginá que el grafo crece a \(|\V| = 10^9\). ¿Sigue
    siendo viable k-hop con el layout actual? ¿Qué cambiarías?
\end{enumerate}
\end{ejercicios}
